% concatenated from Submission_version_for_arXiv.tex
\documentclass[10pt,leqno]{amsart}
\usepackage{graphicx}
\baselineskip=16pt

\usepackage{indentfirst,csquotes}
\usepackage{mathtools}

\topmargin= .5cm
\textheight= 20cm
\textwidth= 32cc
\baselineskip=16pt

\evensidemargin= .9cm
\oddsidemargin= .9cm

\usepackage{amssymb,amsthm,amsmath}
\usepackage{xcolor,paralist,hyperref,titlesec,fancyhdr,etoolbox}

\newtheorem{theorem}{Theorem}[section]
\newtheorem*{theorem*}{Theorem}
\newtheorem{lemma}[theorem]{Lemma}
\newtheorem*{corollary*}{Corollary}
\newtheorem{corollary}[theorem]{Corollary}
\newtheorem*{proposition*}{Proposition}
\newtheorem{proposition}[theorem]{Proposition}
\newtheorem*{acknowledgments*}{Acknowledgments}


\theoremstyle{definition}
\newtheorem{definition}[theorem]{Definition}
\newtheorem{example}[theorem]{Example}
\newtheorem{xca}[theorem]{Exercise}

\theoremstyle{remark}
\newtheorem{remark}[theorem]{Remark}

\numberwithin{equation}{section}

\def\nn{\mathbb{N}}
\def\zz{\mathbb{Z}}
\def\qq{\mathbb{Q}}
\def\rr{\mathbb{R}}
\def\cc{\mathbb{C}}
\def\ff{\mathbb{F}}
\def\pp{\mathbb{P}}
\def\hh{{\mathbb H}}
\def\gcd{{\rm gcd}}
\def\lcm{{\rm lcm}}
\def\ker{{\rm ker}}
\def\rank{{\rm rank}}
\def\Span{{\rm Span}}
\def\im{{\rm Im}}
\def\re{{\rm Re}}
\def\disp{\displaystyle}
\def\supp{{\rm supp}}
\def\th{{\rm th}}
\def\End{{\rm End}}
\def\GL{{\rm GL}}
\def\SL{{\rm SL}}
\def\tr{{\rm tr}}
\def\id{{\rm id}}
\def\op{{\rm op}}
\def\rad{{\rm Rad}}
\def\Inn{{\rm Inn}}
\def\Out{{\rm Out}}
\def\ord{{\rm ord}}
\def\ann{{\rm ann}}
\def\uu{\mathcal{U}}
\def\ss{\mathcal{S}}
\def\mbfzz{\mathbf{Z}}
\def\hom{{\rm Hom}}
\def\mbfee{\mathbf{E}}
\def\Irr{{\rm Irr}}
\def\Cl{{\rm Cl}}
\def\diag{{\rm diag}}
\def\Gal{{\rm Gal}}
\def\Syl{{\rm Syl}}
\def\Inn{{\rm Inn}}
\def\vv{{\rm V}}
\def\uu{{\rm U}}
\def\conj{{\rm conj}}
\def\cen{{\rm Z}}
\def\Fit{{\rm Fit}}
\def\Comp{{\rm Comp}}
\def\aug{{\rm aug}}
\def\O{{\mathcal{O}}}
\def\a{{\alpha}}
\def\b{{\beta}}
\def\p{{\mathfrak{p}}}
\def\gl{{\rm GL}}
\newcommand{\gal}{\operatorname{Gal}}
\newcommand{\aut}{\operatorname{Aut}}
\newcommand{\cl}{\operatorname{Cl}}
\newcommand{\sgn}{\operatorname{sgn}}
\newcommand{\cha}{\operatorname{char}}
\newcommand{\Mod}{\operatorname{mod}}
\newcommand{\Tr}{\operatorname{Tr}}
\newcommand{\res}{\operatorname{Res}}

\titleformat{\section}
{\normalfont\Large\bfseries\centering}
{\thesection.}
{1em}
{}

\hypersetup{ colorlinks=true, linkcolor=black, filecolor=black, urlcolor=black }
\def\proof{\noindent {\it Proof. $\, $}}
\def\endproof{\hfill $\Box$ \vskip 5 pt }

\usepackage{lipsum}

\begin{document}
\title{Resultants of dynatomic polynomials of $x^d$} %%%%%%%%%%%%
\author{Chih-Chiang Kao}
\date{}
\address{Department of Mathematics, Institute of Science Tokyo, 2-12-1 Ookayama, Meguro-ku, 152-8550, Tokyo, Japan}
\email{kao.c.aa@m.titech.ac.jp}
\maketitle

\let\thefootnote\relax
\footnotetext{MSC2020: Primary 13P15, Secondary 11B83 .} %%%%%%%%%%
\footnotetext{Key words and phrases. Resultant, dynatomic polynomial, cyclotomic polynomial.}

\begin{abstract}
Let $K$ be a field of characteristic zero, and let $\phi(x)\in K[x]$ be a polynomial of degree at least 2. Denote the $n$-th iterate of $\phi$ by $\phi^n$. The $n$-th dynatomic polynomial of $\phi$ is defined by $$\Phi_{\phi,n}(x) \coloneqq \prod_{k\mid n}(\phi^k(x)-x)^{\mu(n/k)}.$$
In this paper, we specialize to the case $\phi(x) = x^d$. We first establish several properties of $\Phi_{\phi,n}$ that are analogous to those of cyclotomic polynomials. We then combine these properties with known results on resultants of cyclotomic polynomials to determine the resultants of dynatomic polynomials. In particular, for $1\leq n< m$, we show that $\res(\Phi_{\phi,n},\Phi_{\phi,m}) = 1$
if and only if $n\nmid m$, and we obtain an explicit formula for $\res(\Phi_{\phi,1},\Phi_{\phi,m})$.
\end{abstract} %%%%%%%%%

\bigskip

\section{Introduction}

Let $K$ be a field of characteristic zero. Let $\phi(x)\in K[x]$ be a polynomial with $\deg\phi\geq 2$. For each positive integer $k$, let $\phi^k(x)$ denote the $k$-th iterate of $\phi$. We also set $\phi^0(x) = x$. We begin by defining the dynatomic polynomial associated with $\phi$.

\begin{definition}
    The $n$-th \textit{dynatomic polynomial} $\Phi_{\phi,n}(x)\in K[x]$ associated with $\phi$ is defined by
    $$\Phi_{\phi,n}(x) = \prod_{k\mid n}(\phi^k(x)-x)^{\mu(n/k)},$$
    where $\mu(\cdot)$ denotes the Möbius function.
\end{definition}

Although the above definition yields a rational function, since the Möbius function may take negative values, Morton and Patel \cite[Theorem 2.5]{Morton94} proved that $\Phi_{\phi,n}(x) = \prod_{k\mid n}(\phi^k(x)-x)^{\mu(n/k)}$ is in fact a polynomial over $K$. For more details on dynatomic polynomials, we refer the reader to \cite[Section 4.1]{Silverman07}.

A point $\alpha\in\pp^1_K$ is called a periodic point of $\phi$ of period $n$ if $\phi^n(\alpha) = \alpha$. That is, $\alpha$ is a root of the polynomial $\phi^n(x)-x$. We say that $\alpha$ has an exact period $n$ if $\phi^n(\alpha) = \alpha$ and $\phi^i(\alpha)\neq \alpha$ for all $1\leq i< n$. We also say that $\alpha$ has a formal period $n$ if $\Phi_{\phi,n}(\alpha) = 0$.

It follows immediately from the definitions that
$$\alpha\text{ has an exact period }n \Rightarrow \alpha\text{ has a formal period }n \Rightarrow \alpha\text{ has period }n.$$
On the other hand, the example $\phi(x) = x^2-\frac{3}{4}$ in $\qq[x]$ shows that neither converse implication holds.

The $n$-th cyclotomic polynomial $\Phi_n(x)$, whose roots are precisely the primitive $n$-th roots of unity, is defined by
$$\Phi_n(x) = \prod_{k\mid n}(x^k-1)^{\mu(n/k)}.$$

Let $\alpha$ be a root of the dynatomic polynomial $\Phi_{\phi,n}(x)$. The multiplier $\omega_n(\alpha)$ is defined as the derivative of $\phi^n$ at $\alpha$. In \cite{Vivaldi92}, it is proved that there exists a monic polynomial $\delta_n(x)\in K[x]$ satisfying
$$(\delta_n(x))^n = \prod_{\Phi_{\phi,n}(\alpha) = 0}(x-\omega_n(\alpha)),$$
which is called the $n$-th multiplier polynomial of $\phi$. In \cite{Morton95}, it is proved that:
for any divisor $m$ of $n$ with $m<n$, we have
$$\res(\Phi_{\phi,n},\Phi_{\phi,m}) = \pm\Delta_{n,m}^m,$$
where $\Delta_{n,m} = \res(\Phi_{n/m}(x),\delta_m(x))$.

The resultant of $\Phi_{\phi,n}$ and $\Phi_{\phi,m}$ can be expressed in terms of the resultant of a cyclotomic polynomial and the multiplier polynomial $\delta_m$. However, it is difficult to determine $\delta_m$ explicitly. 

In this paper, we restrict our attention to the case $\phi(x) = x^d$ and use known results on resultants of cyclotomic polynomials to prove the following theorem.
\begin{theorem*}
    For $1\leq n<m$, we have 
    \begin{enumerate}
        \item $\res(\Phi_{\phi,1},\Phi_{\phi,m}) = \Phi_m(d)^{d-1}$,
        \item $\res(\Phi_{\phi,n},\Phi_{\phi,m}) = 1$ if and only if $n\nmid m$.
    \end{enumerate}
\end{theorem*}
We also describe $\res(\Phi_{\phi,n},\Phi_{\phi,m})$ when $n\mid m$. 

In the next section, we review some well-known results on resultants that will be used throughout the remainder of the paper. In Section 3, we establish some properties of $\Phi_{d,n}$ analogous to those of cyclotomic polynomials. In the final section, we determine the resultant $\res(\Phi_{\phi,1},\Phi_{\phi,m})$ for $m>1$ and then investigate $\res(\Phi_{\phi,n},\Phi_{\phi,m})$ for $1\leq n<m$.

\begin{acknowledgments*}
    The author expresses sincere gratitude to Professor Yuichiro Taguchi for his continued help and suggestions.
\end{acknowledgments*}




\section{Facts of resultants}

In this section, we briefly review the definition and some well-known properties of resultants that will be used in the following two sections. For further details, we refer the reader to \cite[Chapter 12]{Gelfand94} or \cite[\textup{Chapter 5}]{Waerden70}.

Given any two polynomials 
$$A = a_nx^n+\cdots+a_1x+a_0\text{ and }B = b_mx^m+\cdots+b_1x+b_0,$$
we define the resultant $\res(A,B)$ of $A$ and $B$ as the determinant of their Sylvester matrix:
$$\begin{array}{c@{\hspace{0.2em}}l}
    \begin{vmatrix}
        a_n & a_{n-1} & \cdots  & a_1    & a_0     &        &        &     \\
            & a_n     & a_{n-1} & \cdots & a_1     & a_0    &        &     \\
            &         & \ddots  & \ddots & \ddots  & \ddots & \ddots &     \\
            &         &         & a_n    & a_{n-1} & \cdots & a_1    & a_0 \\
        b_m & b_{m-1} & \cdots  & b_1    & b_0     &        &        &     \\
            & b_m     & b_{m-1} & \cdots & b_1     & b_0    &        &     \\
            &         & \ddots  & \ddots & \ddots  & \ddots & \ddots &     \\
            &         &         & b_m    & b_{m-1} & \cdots & b_1    & b_0
    \end{vmatrix} & 
    \begin{array}{l}
        \left. \rule{0pt}{7ex} \right\} m \text{ rows} \\[3.5ex]
        \left. \rule{0pt}{7ex} \right\} n \text{ rows}
    \end{array}
\end{array},$$
where the remaining entries are all zero.

\begin{proposition} 
    Suppose that
    $$A = a_n(x-\alpha_1)\cdots(x-\alpha_n)\text{ and }B = b_m(x-\beta_1)\cdots(x-\beta_m).$$
    Then:
    \begin{enumerate}
        \item $\res(A,B) = a_n^m\prod_{i = 1}^nB(\alpha_i) = (-1)^{mn}b_m^n\prod_{j = 1}^mA(\beta_j).$
        \item $\res(B,A) = (-1)^{mn}\res(A,B)$.
        \item If $C$ is another polynomial, then $\res(AB,C) = \res(A,C)\res(B,C)$.
    \end{enumerate}
\end{proposition}

To compute the resultants of dynatomic polynomials of $x^d$, it is helpful to know the resultants of cyclotomic polynomials. Lehmer \cite{Lehmer30}, Diederichsen \cite{Diederichsen40}, Apostol \cite{Apostol70}, Louboutin \cite{Louboutin97}, and Dresden \cite{Dresden12} have provided different proofs for the resultants of cyclotomic polynomials. We omit the proofs and state the result below.

\begin{theorem}
    For integers $0<n<m$, the resultant of $\Phi_n$ and $\Phi_m$ is
    $$\res(\Phi_n,\Phi_m) = \begin{cases}
        p^{\varphi(n)} &\text{ if }m/n\text{ is a power of a prime }p,\\
        1 &\text{otherwise},
    \end{cases}$$
    where $\varphi(\cdot)$ denotes Euler's totient function.
\end{theorem}

Note that $\deg\Phi_n = \varphi(n)$ is even for $n\geq 3$. Hence, for integers $0<m<n$ with $n\geq 3$, Proposition 2.1(2) gives $\res(\Phi_m,\Phi_n) = \res(\Phi_n,\Phi_m)$. On the other hand, $\res(\Phi_2,\Phi_1) = -\res(\Phi_1,\Phi_2) = -2$.

\section{Basic properties of dynatomic polynomials for $x^d$}

In this section, we study the dynatomic polynomials of $\phi(x) = x^d$, where $d\geq 2$. We first recall the definition of the dynatomic polynomial of $\phi(x)$:
$$\Phi_{\phi,n}(x) = \prod_{k\mid n}(\phi^k(x)-x)^{\mu(n/k)} = \prod_{k\mid n}(x^{d^k}-x)^{\mu(n/k)},$$
where $\mu(\cdot)$ is the Möbius function. For convenience, we denote $\Phi_{\phi,n}$ by $\Phi_{d,n}$.

We recall the following properties of Euler's totient function and the Möbius function. For every positive integer $n$, we have
$$\sum_{k\mid n}\varphi(k) = n,$$
and
$$\sum_{k\mid n}\mu(k) = \begin{cases}
    1&\text{ if }n = 1,\\
    0&\text{ if }n>1.
\end{cases}$$
We also recall the following multiplicative version of the Möbius inversion formula. Let $f$ and $g$ be arithmetic functions. Then, for every positive integer $n$, the following equivalent holds:
$$g(n) = \prod_{k\mid n}f(k)\Leftrightarrow f(n) = \prod_{k\mid n}g(k)^{\mu(n/k)}.$$
We refer the reader to \cite[Chapter 2]{Apostol76} for proofs of these results.

In Section 5.1 of \cite{Vivaldi92}, a decomposition of $\Phi_{2,n}$ is given for $n\geq 1$. In this section, we prove that $\Phi_{d,n}$ factors into a product of cyclotomic polynomials for all $d,n\geq 2$.

\begin{lemma}
    For every $d,n\geq 2$,
    $$\Phi_{d,n}(x) = \prod_{\ord_k(d) = n}\Phi_k(x),$$
    where $\ord_k(d)$ denotes the multiplicative order of $d$ modulo $k$.
\end{lemma}

\begin{proof}
    Note that
    $$\Phi_{d,n}(x) = \prod_{m\mid n}(x^{d^m}-x)^{\mu(n/m)} \text{ and }x^{d^m-1}-1 = \prod_{k\mid d^m-1}\Phi_k(x).$$
    Hence, we obtain
    \begin{align*}
        \Phi_{d,n} &= \prod_{m\mid n}\left(x\prod_{k\mid d^m-1}\Phi_k\right)^{\mu(n/m)}\\
        &= x^{\sum_{m\mid n}\mu(n/m)}\cdot\prod_{m\mid n}\prod_{k\mid d^m-1}\Phi_k^{\mu(n/m)}\\
        &= \prod_{m\mid n}\prod_{k\mid d^m-1}\Phi_k^{\mu(n/m)}
    \end{align*}
    since $\sum_{m\mid n}\mu(n/m) = 0$ when $n\geq 2$. 
    
    Fix an integer $n$ and a divisor $m$ of $n$. Then $k\mid d^m-1$ implies that $k\mid d^n-1$. Conversely, fix a divisor $k$ of $d^n-1$, we wish to determine the divisors $m$ such that $k\mid d^m-1$. Note that $k\mid d^m-1$ if and only if $\ord_k(d)\mid m$. Letting $m = \ord_k(d)\cdot t$, the divisor $m$ divides $n$ if and only if $t$ divides $n/\ord_k(d)$. Hence, we obtain
    $$\Phi_{d,n} = \prod_{m\mid n}\prod_{k\mid d^m-1}\Phi_k^{\mu(n/m)} = \prod_{k\mid d^n-1}\prod_{t\mid \frac{n}{\ord_k(d)}}\Phi_k^{\mu(\frac{n}{\ord_k(d)\cdot t})}.$$
    Using the property of the Möbius function, we obtain
    $$\prod_{t\mid \frac{n}{\ord_k(d)}}\Phi_k^{\mu(\frac{n}{\ord_k(d)\cdot t})} = \begin{cases}
        \Phi_k &\text{ if }n = \ord_k(d)\\
        1&\text{ if }n\neq \ord_k(d).
    \end{cases}$$
    Since $\ord_k(d) = n$ implies that $k\mid d^n-1$, we conclude that $\Phi_{d,n} = \prod_{\ord_k(d) = n}\Phi_k$.
\end{proof}

\begin{remark}
    Morton and Patel \cite[Theorem 2.5]{Morton94} proved that if $\phi(x)\in K[x]$ with $\deg\phi(x)\geq 2$, then 
    $$\Phi_{\phi,n}(x) = \prod_{k\mid n}(\phi^k(x)-x)^{\mu(n/k)}$$
    is also a polynomial. In the special case $\phi(x) = x^d$, Lemma 3.1 yields an independent proof of this fact by expressing $\Phi_{d,n}$ as a product of cyclotomic polynomials.
\end{remark}

Since $\Phi_n(x) = \prod_{k\mid n}(x^k-1)^{\mu(n/k)}$, the dynatomic polynomials $\Phi_{d,n}$ may be viewed as dynamical analogues of cyclotomic polynomials. The following proposition gives dynamical analogues of well-known properties of cyclotomic polynomials.

\begin{proposition}For integers $d,n\geq 2$, the following hold.
    \begin{enumerate}
        \item $\deg\Phi_{d,n}(x) = \sum_{k\mid n}\mu(n/k)d^k = \sum_{\ord_k(d) = n}\varphi(k)$.
        \item $x^{d^n}-x = \prod_{k\mid n}\Phi_{d,k}(x)$.
        \item $\Phi_{d,n}(x)$ is a palindromic polynomial of even degree.
    \end{enumerate}
\end{proposition}

\begin{proof}~
    \begin{enumerate}
        \item By definition, $\Phi_{d,n}(x) = \prod_{k\mid n}(x^{d^k}-x)^{\mu(n/k)}$. Since $d^k>1$ for every $k$, the first equality follows immediately. For the second equality, since the degree of the $n$-th cyclotomic polynomial is $\varphi(n)$ and $\Phi_{d,n}(x) = \prod_{\ord_k(d) = n}\Phi_k(x)$, we obtain $\deg \Phi_{d,n}(x) = \sum_{\ord_k(d) = n}\varphi(k)$.
        \item Since $\Phi_{d,n}(x) = \prod_{k\mid n}(x^{d^k}-x)^{\mu(n/k)}$, the result follows from the multiplicative version of the Möbius inversion formula.
        \item Note that $\deg\Phi_{d,n}(x) = \sum_{k\mid n}\mu(n/k)d^k$. If $d$ is even, it is clear that $\deg\Phi_{d,n}$ is even. If $d$ is odd, then
        $$\sum_{k\mid n}\mu(n/k)d^k\equiv \sum_{k\mid n}\mu(n/k) = 0\pmod 2.$$
        Therefore, the degree of $\Phi_{d,n}$ is even.

        Since $\Phi_{d,n}$ is a product of cyclotomic polynomials, and each cyclotomic polynomial is palindromic, it follows that $\Phi_{d,n}(x)$ is also palindromic.
    \end{enumerate}
\end{proof}

For every $n\geq 1$, the cyclotomic polynomial $\Phi_n$ is irreducible over $\qq$. In contrast, for integers $d,n\geq 2$, the polynomial $\Phi_{d,n}$ is not always irreducible over $\qq$. We provide the following characterization of the irreducibility of $\Phi_{d,n}$. Before proving this proposition, we recall the following theorem of Zsigmondy \cite{Zsigmondy92}.

\begin{theorem}
    Let $a>b>0$ be coprime integers, and let $n\geq 1$. Then, there exists a prime divisor of $a^n-b^n$ that does not divide $a^k-b^k$ for any $1\leq k<n$, with the following exceptions:
    \begin{enumerate}
        \item $n = 1$, and $a-b = 1$,
        \item $n = 2$, and $a+b$ is a power of $2$,
        \item $(a,b,n) = (2,1,6)$.
    \end{enumerate}
\end{theorem}

\begin{proposition}
    $\Phi_{d,n}(x)$ is irreducible over $\qq$ if and only if $d^n-1$ is a prime number.
\end{proposition}

\begin{proof}
    Note that $\Phi_{d,n} = \prod_{\ord_k(d) = n}\Phi_k(x)$. If $d^n-1$ is a prime number, then $\Phi_{d,n}(x) = \Phi_{d^n-1}(x)$, which is irreducible over $\qq$. In particular, $d^n-1$ can be prime only if $d = 2$ and $n$ is a prime number, that is, $d^n-1$ is a Mersenne prime.

    Conversely, suppose that $d^n-1$ is not a prime number. If the triple $(d,1,n)$ is not one of the exceptional cases in Zsigmondy's theorem, then there exists a prime $p\mid d^n-1$ and $p\nmid d^k-1$ for all $1\leq k<n$. Thus, $\ord_p(d) = n$. Since $\ord_{d^n-1}(d) = n$, it follows that $\Phi_p\cdot\Phi_{d^n-1}\mid \Phi_{d,n}$.  If $(d,1,n) = (2,1,6)$, then $\Phi_{2,6} = \Phi_{63}\cdot\Phi_{21}\cdot\Phi_9$ which is reducible over $\qq$. If $n = 2$ and $d+1 = 2^c$ for some $c\geq 2$, then $d-1 = 2^c-2$ and so $d^2-1 = 2^{c+1}(2^{c-1}-1)$. Hence, $4\mid d^2-1$ and $4\nmid d-1$, which implies that $\ord_4(d) = 2$. Thus, $\Phi_4\cdot\Phi_{d^2-1}\mid \Phi_{d,n}$.
\end{proof}

In the remainder of this section, we compute the values of $\Phi_{d,n}(x)$ at $x = 0$ and $x = 1$. Since $\Phi_{d,n}$ is a product of cyclotomic polynomials, we first recall the values of cyclotomic polynomials at $x = 0$ and $x = 1$. For $n\geq 2$, it is well known that $\Phi_n(0) = 1$. On the other hand,
\begin{align*}
    \Phi_n(1) = \begin{cases}
        1&\text{ if }n\text{ is not a prime power},\\
        p&\text{ if }n = p^a\text{ is a prime power with }a\geq 1.
    \end{cases}
\end{align*}
We refer the reader to \cite[Section 4.1]{Lang94} for a proof.

Note that $\Phi_{d,n}(1) = \prod_{\ord_k(d) = n}\Phi_k(1)$. Thus, for every prime divisor $p$ of $\Phi_{d,n}(1)$, computing $v_p(\Phi_{d,n}(1))$ amounts to determining the number of positive integers $a$ such that $\ord_{p^a}(d) = n$. Define
$$S(p,d,n)\coloneqq\{a\geq 1:\ord_{p^a}(d) = n\}.$$
We prove that $v_p(\Phi_{d,n}(1)) = |S(p,d,n)| = v_p(\Phi_n(d))$.

We first recall the lifting-the-exponent lemma. We refer the reader to \cite{Par11} for further details. For completeness, we include a proof. Let $p$ be a prime number, and let $n$ be a positive integer, we denote the valuation of $n$ at $p$ by $v_p(n)$.

\begin{lemma}[lifting-the-exponent lemma]
    Let $p$ be a prime number, and let $x,y$ be integers such that $p\nmid x$ and $p\nmid y$.
    \begin{enumerate}
        \item If $p\geq 3$ and $x\equiv y\pmod p$, then
        $$v_p(x^k-y^k) = v_p(x-y)+v_p(k).$$
        \item If $p = 2$ and $x\equiv y\pmod 4$, then 
        $$v_2(x^k-y^k) = v_2(x-y)+v_2(k).$$
    \end{enumerate}
\end{lemma}

\begin{proof}
    For any prime number $p$, we first show that $v_p(x^k-y^k) = v_p(x-y)$ when $\gcd(p,k) = 1$. Note that
    $$x^k-y^k = (x-y)(x^{k-1}+x^{k-2}y+\cdots+xy^{k-2}+y^{k-1}),$$
    it is sufficient to show that $p$ does not divide $x^{k-1}+x^{k-2}y+\cdots+xy^{k-2}+y^{k-1}$. Since $x\equiv y\pmod p$, $p\nmid xy$, and $\gcd(p,k) = 1$, we have
    $$x^{k-1}+x^{k-2}y+\cdots+xy^{k-2}+y^{k-1}\equiv kx^{k-1}\not\equiv 0\pmod p.$$
    Hence, we can assume that $k$ is a power of $p$. We now consider separately the cases where $p$ is an odd prime and $p = 2$.
    \begin{enumerate}
        \item Suppose that $k = p$ is an odd prime number. Then 
        $$x^{p-1}+x^{p-2}y+\cdots+xy^{p-2}+y^{p-1}\equiv px^{p-1}\equiv 0\pmod p,$$
        which implies that $v_p(x^k-y^k) = v_p(x-y)+1$. Suppose now that $k = p^v$. Then $p\nmid x^{p^v}$, $p\nmid y^{p^v}$, $x^{p^v}\equiv y^{p^v}\pmod p$, we obtain
        $$v_p(x^{p^v}-y^{p^v}) = v_p(x^{p^{v-1}}-y^{p^{v-1}})+1 = \cdots = v_p(x-y)+v.$$
        \item Let $p = 2$. Since $2\nmid x$ and $2\nmid y$, it follows that $x^{2^i}\equiv y^{2^i}\equiv 1\pmod 4$ for all positive integers $i$. Moreover, since $4\mid x-y$, we have $x+y\equiv 2\pmod 4$. Thus, 
        $$v_2(x^{2^v}-y^{2^v}) = v_2(x-y)+\sum_{i = 0}^{v-1} v_2(x^{2^{i}}+y^{2^i}) =  v_2(x-y)+v.$$
    \end{enumerate}
\end{proof}

\begin{lemma}
    Let $d,n\geq 2$ be positive integers. Let $p$ be an odd prime divisor of $\Phi_n(d)$. If $p\nmid n$, then $\ord_p(d) = n$, and $v_p(\Phi_n(d)) = v_p(d^n-1)$. If $p\mid n$ and $n = n'\cdot p^k$ where $k = v_p(n)$, $\gcd(n',p) = 1$. Then $\ord_p(d) = n'$ and $v_p(\Phi_n(d)) = 1$.

    For $p = 2$, one has
    $$v_2(\Phi_n(d)) = \begin{cases}
        v_2(d+1),&\text{ if }n = 2,\\
        1,&\text{ if }n>2\text{ is a power of }2,\text{ and }d\text{ is odd},\\
        0,&\text{ otherwise.}
    \end{cases}$$
\end{lemma}

\begin{proof}
    We first consider the case where $p\nmid n$. Since $p\mid \Phi_n(d)$ and $\Phi_n(d)\mid d^n-1$, it follows that $p\mid d^n-1$. Let $m = \ord_p(d)$. Then $m\mid n$. Moreover, $d^m-1 = \prod_{t\mid m}\Phi_t(d)$ and hence $p\mid \Phi_t(d)$ for some divisor $t$ of $m$. If $m<n$, then $t<n$. Consider $x^n-1$ in $\ff_p[x]$. Since $p\mid\Phi_n(d)$ and $p\mid \Phi_t(d)$, the element $d$ is a common root of $\Phi_n(x)$ and $\Phi_t(x)$. Hence, $d$ is a multiple root of $x^n-1$. However since $p\nmid n$, we have $\gcd(x^n-1,nx^{n-1}) = 1$, contradicting that $x^n-1$ has a multiple root. Thus, $m = n$. Consequently, for every proper divisor $r$ of $n$, we have $p\nmid\Phi_r(d)$. It follows that 
    $$v_p(d^n-1) = \sum_{r\mid n}v_p(\Phi_r(d)) = v_p(\Phi_n(d)).$$

    Next, suppose $p\mid n$, and $n = n'\cdot p^k$ where $k = v_p(n), \gcd(n',p) = 1$. Note that
    $$\Phi_n(x) = \Phi_{n'\cdot p^k}(x) = \Phi_{n'\cdot p}(x^{p^{k-1}}) = \Phi_{n'}(x^{p^k})/\Phi_{n'}(x^{p^{k-1}}).$$
    Hence, $\Phi_n(d) = \Phi_{n'}(d^{p^k})/\Phi_{n'}(d^{p^{k-1}})$. If $p\mid \Phi_n(d)$, then $p$ also divides $\Phi_{n'}(d^{p^k})$. Thus,
    $$\Phi_{n'}(d^{p^k})\equiv \Phi_{n'}(d^{p^{k-1}})\equiv \cdots\equiv \Phi_{n'}(d)\equiv 0\pmod p.$$
    Applying the result proved above, since $p\nmid n'$ and $p\mid \Phi_{n'}(d)$, it follows that $\ord_p(d) = n'$. Since $\gcd(n',p) = 1$ and $p\mid \Phi_{n'}(d)$, we can deduce that $p\mid d^{n'}-1$. According to Lemma 3.6, $v_p(d^{n'p}-1) = v_p(d^{n'}-1)+1$, which implies that $p\mid d^{n'p}-1$. Applying the same argument repeatedly, we obtain $v_{p}(d^{n'p^2}-1) = v_p(d^{n'p}-1)+1$. By induction, $v_p(d^{n'p^i}-1) = v_p(d^{n'p^{i-1}})+1$ for all $i\geq 1$. Hence,
    $$v_p(\Phi_n(d)) = v_p(\Phi_{n'}(d^{p^k}))-v_p(\Phi_{n'}(d^{p^{k-1}})) = v_p(d^{n'p^k}-1)-v_p(d^{n'p^{k-1}}-1) = 1.$$

    Finally, let $p = 2$. If $n = 2$, then $\Phi_n(d) = d+1$ and so $v_2(\Phi_n(d)) = v_2(d+1)$. If $n>2$ is a power of 2, then $\Phi_n(d) = d^{n/2}+1$. Hence,
    $$v_2(\Phi_n(d)) = \begin{cases}
        1&\text{ if }d\text{ is odd},\\
        0&\text{ if }d\text{ is even}.\\
    \end{cases}$$
    For the remaining cases, if $d$ is even, then $\Phi_n(d)\equiv \Phi_n(0)\equiv 1\pmod 2$. If $d$ is odd, then $\Phi_n(d)\equiv \Phi_n(1)\equiv 1\pmod 2$. In all cases, $v_2(\Phi_n(d)) = 0$.
\end{proof}
\vspace{1em}
\begin{lemma}~
    \begin{enumerate}
        \item Let $p$ be a prime and let $d$ be a positive integer with $p\nmid d$. Then
        $$\ord_{p^a}(d) = \min\{r\in\nn\mid v_p(d^r-1)\geq a\}.$$
        \item Let $p$ be an odd prime number. If $\ord_p(d) = n$ and $v_p(d^n-1) = k$, then $\ord_{p^a}(d) = n\cdot p^{\max(0,a-k)}$.
        \item Let $p = 2$. If $d\equiv 1\pmod 4$ and $v_2(d-1) = k$, then $\ord_{2^a}(d) = 2^{\max(0,a-k)}$.
        If $d\equiv 3\pmod 4$ and $v_2(d+1) = k$, then
        $$\ord_{2^a}(d) = \begin{cases}
            1 &\text{ if }a = 1,\\
            2^{\max(1,a-k)}&\text{ if }a\geq 2.
        \end{cases}$$
    \end{enumerate}
\end{lemma}

\begin{proof}~
    \begin{enumerate}
        \item Let $\ord_{p^a}(d) = n$. Then $d^n\equiv 1\pmod {p^a}$ and $d^i\not\equiv 1\pmod {p^a}$ for all $0<i<n$. On the other hand, denote $\min\{r\in\nn\mid v_p(d^r-1)\geq a\}$ by $m$. Then $v_p(d^m-1)\geq a$ and $v_p(d^i-1)<a$ for all $0< i<m$.

        Since $v_p(d^m-1)\geq a$ implies that $d^m\equiv 1\pmod{p^a}$, we have $n\leq m$. Conversely, $d^n\equiv 1\pmod{p^a}$ implies that $v_{p}(d^n-1)\geq a$, and hence $m\leq n$. Therefore, $n = m$.
        \item Let $n' = \ord_{p^a}(d)$. Since $d^{n'}\equiv 1\pmod p$, we have $n\mid n'$. Write $n' = nx$ for some $x\geq 1$. By Lemma 3.6, 
        $$v_p(d^{nx}-1) = v_p(d^n-1)+v_p(x) = k+v_p(x).$$
        By (1), $n'$ is the smallest positive integer $nx$ such that $k+v_p(x)\geq a$. Thus, the smallest possible value of $x$ is $x = p^{\max(0,a-k)}$, and hence
        $$\ord_{p^a}(d) = n' = n\cdot p^{\max(0,a-k)}.$$
        \item Suppose first that $d\equiv 1\pmod 4$. Then $k = v_2(d-1)\geq 2$. By Lemma 3.6,
        $$v_2(d^r-1) = v_2(d-1)+v_2(r)$$
        for every $r\geq 1$. Hence, by (1), the order of $d$ modulo $2^a$ is determined by the smallest $r$ satisfying $k+v_2(r)\geq a$. Thus,
        $$\ord_{2^a}(d) = 2^{\max(0,a-k)}.$$

        Now suppose that $d\equiv 3\pmod 4$. Then $\ord_2(d) = 1$ and $\ord_4(d) = 2$. Moreover, if $r$ is odd, then $v_2(d^r-1) = v_2(d-1) = 1$, so for $a\geq 2$ the order must be even. Write $r = 2s$. Lemma 3.6 gives
        $$v_2(d^{2s}-1) = v_2(d^2-1)+v_2(s) = 1+k+v_2(s).$$
        By (1), for $a\geq 2$ the smallest such $s$ satisfies $1+k+v_2(s)\geq a$. Hence $s = 2^{\max(0,a-k-1)}$, and therefore $\ord_{2^a}(d) = 2s = 2^{\max(1,a-k)}$.
    \end{enumerate}
\end{proof}

\begin{lemma}
    Let $p$ be a prime number, and let $d,n\geq 2$ be integers satisfying $\gcd(p,d) = 1$. Recall that
    $$S(p,d,n)\coloneqq\{a\geq 1:\ord_{p^a}(d) = n\}.$$
    Then 
    $$v_p(\Phi_n(d)) = |S(p,d,n)|.$$
\end{lemma}

\begin{proof}
    By Lemma 3.7, it suffices to show the following:
    \begin{enumerate}
        \item If $p\nmid \Phi_n(d)$, then there does not exist $a\geq 1$ such that $p^a\mid d^n-1$ and $\ord_{p^a}(d) = n$.
        \item If $p$ is odd, $p\mid \Phi_n(d)$, and $p\nmid n$, then $|S(p,d,n)| = v_p(d^n-1)$.
        \item If $p$ is odd, $p\mid \Phi_n(d)$, and $p\mid n$, then $|S(p,d,n)| = 1$.
        \item If $2\mid\Phi_n(d)$, then 
        $$|S(p,d,n)| = \begin{cases}
            v_2(d+1) &\text{ if }n = 2,\\
            1&\text{ if } n>2\text{ is a power of 2 and }d\text{ is odd},\\
            0&\text{ otherwise}.
        \end{cases}$$
    \end{enumerate}
    Assume first that $p$ is odd. Let $v_p(d^n-1) = k$. Then $a\in S(p,d,n)$ implies that $1\leq a\leq k$ and so $\ord_{p^a}(d) = \ord_p(d)$ by Lemma 3.8.
    \begin{enumerate}
        \item Suppose $p\nmid \Phi_n(d)$, and there exists $a\geq 1$ such that $p^a\mid d^n-1$ and $\ord_{p^a}(d) = n$. Then for all $e<n$, one has $d^e\not\equiv 1\pmod {p^a}$. Since $d^n-1 = \Phi_n(d)\cdot\prod_{\substack{e\mid n\\e<n}}\Phi_e(d)$, it follows that there exists $e_0<n$ and $e_0\mid n$ such that $p\mid \Phi_{e_0}(d)$, and hence $p\mid d^{e_0}-1$. So $\ord_{p^a}(d) = \ord_p(d)$ divides $e_0$ which is a contradiction. Hence, $S(p,d,n) = \emptyset$.
        \item If $p\mid \Phi_n(d)$ and $p\nmid n$, then $\ord_{p^a}(d) = \ord_p(d) = n$ for all $1\leq a\leq k$. Therefore, $|S(p,d,n)| = k = v_p(d^n-1)$.
        \item Suppose $p\mid \Phi_n(d)$ and $p\mid n$. If we write $n = n'\cdot p^r$ where $r = v_p(n)$, then $\ord_p(d) = n'$. Let $v_p(d^{n'}-1) = k'$, then $\ord_{p^a}(d) = n'\cdot p^{\max(0,a-k')}$. Hence
        $$\ord_{p^a}(d) = n\Leftrightarrow a-k' = r \Leftrightarrow a = k'+r,$$
        which implies $|S(p,d,n)| = 1 = v_p(\Phi_n(d))$.
        \item For $n = 2$, we have $v_2(\Phi_n(d)) = v_2(d+1)$. If $d\equiv 1\pmod 4$, we denote $v_2(d-1)$ by $k$. Then $k\geq 2$ and $\ord_{2^a}(d) = 2^{\max(0,a-k)}$. Suppose $a\in S(p,d,n)$, then $2^{\max(0,a-k)} = \ord_{2^a}(d) = 2$ if and only if $a = k+1$. So $|S(p,d,n)| = 1 = v_2(d+1) = v_2(\Phi_n(d))$. Suppose $d\equiv 3\pmod 4$, and $v_2(d+1) = k'\geq 2$. If $a = 1$, then $\ord_2(d) = 1$ and so $1\not\in S(p,d,n)$. If $a\geq 2$, then $2^{\max(1,a-k')} = \ord_{2^a}(d) = 2$ if and only if $a\leq 1+k'$. So $|S(p,d,n)| = k' = v_2(d+1) = v_2(\Phi_n(d))$.

        For $n = 2^c$ for some $c> 1$. If $d\equiv 1\pmod 4$, then $\ord_{2^a}(d) = 2^{\max(0,a-k)}$. Suppose $a\in S(p,d,n)$, then $2^{\max(0,a-k)} = \ord_{2^a}(d) = 2^c$ if and only if $a = k+c$. So $|S(p,d,n)| = 1 = v_2(\Phi_n(d))$. Suppose $d\equiv 3\pmod 4$, and $v_2(d+1) = k'\geq 2$. If $a = 1$, then $\ord_2(d) = 1$ and so $1\not\in S(p,d,n)$. If $a\geq 2$, then $2^{\max(1,a-k')} = \ord_{2^a}(d) = 2^c$ if and only if $a\leq c+k'$. So $|S(p,d,n)| = 1 = v_2(\Phi_n(d))$.

        If $n$ is not a power of $2$, or $d$ is even, then $S(p,d,n) = \emptyset$ because $\ord_{2^a}(d)$ is a power of $2$. Thus, $|S(p,d,n)| = 0 = v_2(\Phi_n(d))$.
    \end{enumerate}
\end{proof}

\begin{corollary}
    For integers $d,n\geq 2$, we have $\Phi_{d,n}(0) = 1$ and $\Phi_{d,n}(1) = \Phi_n(d)$.
\end{corollary}

\begin{proof}
    By Lemma 3.1, we have $\Phi_{d,n}(x) = \prod_{\ord_k(d) = n}\Phi_k(x)$. Note that $\Phi_n(0) = 1$ for all $n\geq 2$. If we substitute $x = 0$, then $\Phi_{d,n}(0) = \prod_{\ord_k(d) = n}\Phi_k(0) = 1$. 

    By Lemma 3.9, we have $v_p(\Phi_n(d)) = |S(p,d,n)|$ for any prime $p$. Moreover, $$\Phi_n(1) = \begin{cases}
        p &\text{ if }n = p^a\text{ for some }a>0,\\
        1 &\text{ otherwise.}
    \end{cases}$$
    Hence,
    $$v_p(\Phi_{d,n}(1)) = \sum_{\ord_{p^a}(d) = n}v_p(\Phi_{p^a}(1)) = |S(p,d,n)| = v_p(\Phi_n(d))$$ 
    for all prime numbers $p$. Therefore, $\Phi_{d,n}(1) = \Phi_n(d)$. 
\end{proof}

\section{The resultants of dynatomic polynomials for $x^d$}

Let $1\leq n<m$ be fixed. By Proposition 3.3(3), the degree of $\Phi_{d,m}$ is even. Thus,
$$\res(\Phi_{d,m},\Phi_{d,n}) = (-1)^{\deg \Phi_{d,m}\deg \Phi_{d,n}}\res(\Phi_{d,n},\Phi_{d,m}) = \res(\Phi_{d,n},\Phi_{d,m}).$$
Therefore, it is sufficient to compute $\res(\Phi_{d,n},\Phi_{d,m})$ for $1\leq n<m$. 

Since $\Phi_{d,n} = \prod_{\ord_a(d) = n}\Phi_a$, Proposition 2.1(3) yields
$$\res(\Phi_{d,n},\Phi_{d,m}) = \prod_{\ord_a(d) = n}\prod_{\ord_b(d) = m}\res(\Phi_a,\Phi_b).$$
Since $\ord_a(d)=n$ and $\ord_b(d) = m$ with $n\neq m$, we have $a\neq b$. Hence, by Theorem 2.2, each factor $\res(\Phi_a,\Phi_b)$ is positive. Therefore, $\res(\Phi_{d,n},\Phi_{d,m})$ is a positive integer. 

To determine when this resultant exceeds 1, recall that $\res(\Phi_a,\Phi_b)>1$ only when either $a/b$ or $b/a$ is a power of a prime $p$. Thus, we first investigate when a divisibility relation between $a$ and $b$ occur.

\begin{lemma}
    Let $d,m,n$ be positive integers with $d>1$ and $n\nmid m$. If $a$ and $b$ are positive integers satisfying
    $$\ord_a(d) = n,\quad \ord_b(d) = m,$$
    then $a\nmid b$.
\end{lemma}

\begin{proof}
    Write $m = nq+r$, where $q,r\in\zz$ and $0<r<n$. If $a\mid b$, then $a\mid d^m-1$. Hence,
    $$1\equiv d^m\equiv d^{nq+r}\equiv d^r\pmod a.$$
    Thus, $\ord_a(d)\leq r<n$, which contradicts $\ord_a(d) = n$.
\end{proof}

The following corollary is an immediate consequence of Lemma 4.1.

\begin{corollary}
    Let $n,m$ be positive integers with $n<m$. If $a$ and $b$ are positive integers satisfying
    $$\ord_a(d) = n,\quad \ord_b(d) = m,$$
    then $b\nmid a$. 
\end{corollary}

\begin{theorem}
    If $m>1$, then $\res(\Phi_{d,1},\Phi_{d,m}) = \Phi_m(d)^{d-1}$
\end{theorem}

\begin{proof}
    Note that $\Phi_{d,1}(x)= x\cdot(x^{d-1}-1) = x\cdot\prod_{\ord_a(d) = 1}\Phi_a(x)$, and $\Phi_{d,m}(x) = \prod_{\ord_b(d) = m}\Phi_b(x)$. By Proposition 2.1(1), we have $\res(x,\Phi_{d,m}) = \Phi_{d,m}(0)$. Hence,
    \begin{align*}
        \res(\Phi_{d,1},\Phi_{d,m}) &= \res(x,\Phi_{d,m})\cdot\prod_{\ord_a(d) = 1}\res(\Phi_a,\Phi_{d,m})\\
        &= \Phi_{d,m}(0)\cdot\prod_{\ord_a(d) = 1}\res(\Phi_a,\Phi_{d,m})\\
        &= \prod_{\ord_a(d) = 1}\prod_{\ord_b(d) = m}\res(\Phi_a,\Phi_b).
    \end{align*}
    Since $m>1$, we have $b\nmid a$ by Corollary 4.2. Thus, $\res(\Phi_a,\Phi_b)>1$ can occur only when $b/a$ is a power of a prime $p$. Fix a divisor $a$ of $d-1$ and a prime number $p$. Write $a = a'p^k$, where $k = v_p(a)$ and $\gcd(a',p) = 1$. We determine the number of positive integers $s$ such that $\ord_{ap^s}(d) = m$. Since $\ord_{a'}(d) = 1$, we have
    $$\ord_{ap^s}(d) = \ord_{a'\cdot p^{k+s}}(d) = \lcm(\ord_{a'}(d),\ord_{p^{k+s}}(d)) = \ord_{p^{k+s}}(d).$$
    By Theorem 2.2, $\res(\Phi_a,\Phi_b) = p^{\varphi(a)}$ if and only if $b = p^s\cdot a$ for some $s\geq 1$. Moreover, such an integer $s$ must satisfy $\ord_b(d) = \ord_{p^{k+s}}(d) = m$. Since $\ord_{p^k}(d) = 1$ and $1<m$, it follows that 
    $$\{s\geq 1:\ord_{p^{k+s}}(d) = m\} = \{t\geq 1:\ord_{p^t}(d) = m\} =  S(p,d,m).$$
    Hence, for a fixed $p$, one has
    \begin{align*}
        v_p(\res(\Phi_{d,1},\Phi_{d,m})) &= \sum_{\ord_a(d) = 1}\varphi(a)\cdot|\{s\geq 1:\ord_{p^{k+s}}(d) = m\}|\\
        &= \sum_{\ord_a(d) = 1}\varphi(a)\cdot|S(p,d,m)|\\
        &= \sum_{\ord_a(d) = 1}\varphi(a)\cdot v_p(\Phi_m(d))&&\text{ by Lemma 3.9}\\
        &= (d-1)v_p(\Phi_m(d)).
    \end{align*}
    Therefore,
    $$\res(\Phi_{d,1},\Phi_{d,m}) = \prod_{p\text{ primes}}p^{(d-1)v_p(\Phi_m(d))} = \Phi_m(d)^{d-1}.$$
\end{proof}

Finally, we investigate the resultant $\res(\Phi_{d,n},\Phi_{d,m})$ for $1<n<m$. 

\begin{theorem}
    Let $1<n<m$. The following statements hold. 
    \begin{enumerate}
        \item $\res(\Phi_{d,n},\Phi_{d,m}) = 1$ if and only if $n\nmid m$.
        \item Let $p$ be a prime divisor of $\Phi_{m}(d)$. If $n\mid m$, then
        $$v_p(\res(\Phi_{d,n},\Phi_{d,m}))\geq v_p(\Phi_m(d))\left(\sum_{k\mid n}\mu\left(\frac{n}{k}\right)d^k\right).$$
    \end{enumerate}
\end{theorem}

\begin{proof}
    \begin{enumerate}
        \item First, we suppose $n\nmid m$. Since $n\nmid m$ and $n<m$, it follows that $a\nmid b$ and $b\nmid a$ by Lemma 4.1 and Corollary 4.2. Therefore, $\res(\Phi_a,\Phi_b) = 1$ for every pair $(a,b)$. Hence, $\res(\Phi_{d,n},\Phi_{d,m}) = 1$.

        The converse direction follows immediately from Theorem 4.4(2). For any prime divisor $p$ of $\Phi_m(d)$, since $\sum_{k\mid n}\mu(n/k)d^k = \deg\Phi_{d,n}$, it follows that
        $$v_p(\res(\Phi_{d,n},\Phi_{d,m})) \geq v_p(\Phi_m(d))\cdot\sum_{k\mid n}\mu(n/k)d^k\geq 1.$$
        Hence, if $n\mid m$, then $\res(\Phi_{d,n},\Phi_{d,m})>1$.
        \item Let $p$ be a prime divisor of $\Phi_m(d)$. Note that
        $$\res(\Phi_{d,n},\Phi_{d,m}) = \prod_{\ord_a(d) = n}\prod_{\ord_b(d) = m}\res(\Phi_a,\Phi_b).$$
        By Theorem 2.2, $\res(\Phi_a,\Phi_b) = p^{\varphi(a)}$ if and only if $b/a$ is a power of $p$. Hence,
        \begin{align}
            v_p(\res(\Phi_{d,n},\Phi_{d,m})) &= \sum_{\ord_{a}(d) = n}\sum_{\ord_b(d) = m}v_p(\res(\Phi_a,\Phi_b))\notag\\
            &= \sum_{\ord_a(d) = n}\varphi(a)\cdot |\{t\geq 1: \ord_{ap^t}(d) = m\}|.
        \end{align}
        Define $T(a,p,d,m) = \{t\geq 1:\ord_{a\cdot p^t}(d) = m\}$. Recall that $S(p,d,m) = \{t\geq 1:\ord_{p^t}(d) = m\}$. We claim that
        $$|T(a,p,d,m)|\geq |S(p,d,m)|$$
        whenever $n\mid m$ and $\ord_a(d) = n$. Since $p\mid \Phi_m(d)$, it follows that $p\mid d^m-1$. We denote $v_p(d^m-1)$ by $r$. For each prime divisor $q$ of $m$, we denote $v_p(d^{m/q}-1)$ by $r_q$. Then $t\in S(p,d,m)$ if and only if $p^t\mid d^m-1$ and $p^t\nmid d^{m/q}-1$ for all prime divisors $q$ of $m$. Hence, $\max_{q\mid m}r_q<t\leq r$. Thus, $S(p,d,m) = \{t\geq 1:\max_{q\mid m}r_q<t\leq r\}$ and so $|S(p,d,m)| = r-\max_{q\mid m}r_q$.

        Write $a = a'p^k$, where $k = v_p(a)$ and $\gcd(a',p) = 1$. Then $a\cdot p^t = a'\cdot p^{t+k}$ and 
        $t\in T(a,p,d,m)$ if and only if $a'\cdot p^{t+k}\mid d^m-1$ and $a'\cdot p^{t+k}\nmid d^{m/q}-1$ for all prime divisors $q$ of $m$. To simplify the notation, define
        $$Q(n,m) = \{\text{prime divisors }q\text{ of }m:n\text{ divides }m/q\}.$$
        Suppose $q\in Q(n,m)$. Since $a'\mid a$, it follows that $a'\mid d^n-1$ and so $a'\mid d^{m/q}-1$. In this case, $t+k\leq r$ and $r_q<t+k$. Hence, $T(a,p,d,m)\supseteq \{t\geq 1:\max_{q\in Q(n,m)}(r_q-k)<t\leq r-k\}$, and
        $$|T(a,p,d,m)|\geq r-k-\max_{q\in Q(n,m)}(r_q-k) = r-\max_{q\in Q(n,m)}r_q.$$
        Since $Q(n,m)$ is a subset of the set of prime divisors $q$ of $m$, one has $\max_{q\in Q(n,m)}r_q\leq \max_{q\mid m}r_q$. Hence, 
        $$|T(a,p,d,m)| \geq r-\max_{q\in Q(n,m)}r_q\geq r-\max_{q\mid m}r_q = |S(p,d,m)|.$$
        Substituting this estimate into the equation (4.1) yields
        \begin{align*}
            v_p(\res(\Phi_{d,n},\Phi_{d,m})) &= \sum_{\ord_a(d) = n}\varphi(a)\cdot|T(a,p,d,m)|\\
            &\geq \sum_{\ord_a(d) = n}\varphi(a)\cdot|S(p,d,m)|\\
            &= v_p(\Phi_{m}(d))\cdot\sum_{\ord_a(d) = n}\varphi(a)&&\text{by Lemma 3.9}\\
            &= v_p(\Phi_m(d))\sum_{k\mid n}\mu(n/k)d^k&&\text{by Proposition 3.3(1)}.
        \end{align*}
    \end{enumerate}
\end{proof}

Recall that if $\alpha$ is a periodic point of a polynomial $\phi(x)$, then having exact period $n$ implies having formal period $n$. Moreover, $\phi(x)$ has a point of formal period $n$ whose exact period is a proper divisor $m$ of $n$ if and only if $\res(\Phi_{\phi,n}(x),\Phi_{\phi,m}(x)) = 0$. By Theorem 4.3 and 4.4, we have shown that $\res(\Phi_{d,n},\Phi_{d,m})\neq 0$ for all distinct positive integers $m,n$. Therefore, we obtain the following corollary.

\begin{corollary}
    If $\alpha$ is a periodic point of $\phi(x) = x^d$ with formal period $n$, then $\alpha$ has exact period $n$.
\end{corollary}

\bibliographystyle{plain}
\bibliography{Reference}


\end{document} 
